\selectlanguage{ngerman}%
\frenchspacing%
\begin{ledgroupsized}[r]{120mm}%
\footnotesize%
\pstart%
\noindent\textbf{Überlieferung:}%
\pend%
\end{ledgroupsized}%
\begin{ledgroupsized}[r]{114mm}%
\footnotesize%
\pstart%
\parindent -6mm%
\makebox[6mm][l]{\textit{L}}%
Aufzeichnung:
LH~XXXVII\;4 Bl.\;33.
Ein Blatt\;4\textsuperscript{o};
Ränder beschnitten, zum Teil unregelmäßig;
Fragment eines \href{\Scan{LH_37_04_033wz}}{Wasserzeichens} am oberen Rand;
Papiererhaltungsmaßnahmen.
Eine ganzseitig beschriebene Seite auf Bl.\;%
\href{\Scan{LH_37_04_033r}}{33\,r\textsuperscript{o}}.
Bl.\;\href{\Scan{LH_37_04_033v}}{33\,v\textsuperscript{o}} ist unbeschrieben.
\newline%
Cc~2 Nr.\;481
\pend%
\end{ledgroupsized}%
\vspace{5mm}%
\begin{ledgroup}%
\footnotesize%
\pstart%
\noindent%
\textbf{Datierungsgründe:}
Die Datierung der Aufzeichnung N.\;\ref{Cc2-0481} beruht auf Eigenschaften des Textträgers sowie auf inhaltlichen Betrachtungen.
Siehe die ausführliche Begründung auf S.\;\pageref{Cc2-0481_dat}.%
\pend%
\end{ledgroup}%
\selectlanguage{latin}%
\frenchspacing%
\vspace{8mm}%
\normalsize%
\pstart%
\noindent%
\bracks{\href{\Scan{LH_37_04_033r}}%
{33\,r\textsuperscript{o}}}
\edlabel{LH_37_04_033_033r_Zitate-1a}%
Experiendum
an pila plumbea%
\indexS{pila plumbea}
longius in vacuo quoque%
\indexS{vacuum}
quam
\edtext{%
pluma,%
\indexS{pluma}
quemadmodum%
}{%
\lemma{pluma,
quemadmodum}%
\Bfootnote{%
\textit{erg.\;L}%
}}
\edtext{%
\bracks{in}%
}{%
\lemma{in}%
\Bfootnote{%
\textit{erg. Hrsg.}%
}}
aere%
\indexS{aer (medium resistens)}
jaciatur.
Ego id eventurum non dubito.%
\edlabel{LH_37_04_033_033r_Zitate-1b}
\edlabel{LH_37_04_033_033r_Zitate-2a}%
Demonstro enim hujus rei necessitatem%
\indexS{necessitas rei}
ex regula
\edtext{%
generali%
\indexS{regula generalis}%
\bracks{:}
quo quid%
}{%
\lemma{generali\;,}%
\Bfootnote{%
\textit{(1)}~quod scilicet %
\textit{(2)}~quo quid%
~\textit{L}%
}}
diutius resistit %
virtuti cuidam continuis supplementis auctae%
\indexS{virtus aucta supplementis continuis}%
\indexS{supplementum continuum}%
\bracks{,}
eo major est effectus%
\indexS{effectus virtutis auctae} 
cum tandem vincitur.%
\edlabel{LH_37_04_033_033r_Zitate-2b}%
\pend%
\pstart%
Idque
\edtext{%
ad corpora perspicua applico.%
\indexS{corpus perspicuum}%
}{%
\lemma{ad}%
\Bfootnote{%
\textit{(1)}~vitrum\indexS{vitrum} applico %
\textit{(2)}~corpora perspicua applico.%
~\textit{L}%
}}
\edlabel{LH_37_04_033_033r_Zitate-3a}%
Arbitror enim
manifestum esse,
corpora densiora difficilius cedere ictui%
\indexS{corpus cedens ictui}%
\indexS{corpus densum}
quam rariora,%
\indexS{corpus rarum}
ac proinde diuturniorem
ac per consequens fortiorem
ictum excipere.%
\indexS{ictus diuturnus}%
\indexS{ictus fortis}%
\indexS{ictus exceptus}%
\edlabel{LH_37_04_033_033r_Zitate-3b}
Nec refert
si dicas
quaedam corpora sua natura%
\indexS{natura corporis}
esse plicabilia,%
\indexS{corpus plicabile}
ut aquam,%
\indexS{aqua ut corpus plicabile}
quae non resistet ictui,
sed in latus cedet.
Responsio facilis est
et experimentis conformis,%
\indexS{responsio conformis experimentis}
aquam celerrimo ictui non cedere,%
\indexS{aqua resistens ictui celeri}
sed perinde
ac si dura esset
\edtext{%
repercutere,%
\indexS{aqua repercutiens}
tempore scilicet%
}{%
\lemma{repercutere,}%
\Bfootnote{%
\textit{(1)}~spatio sc %
\textit{(2)}~tempore scilicet%
~\textit{L}%
}}
cedendi negato.%
\indexS{tempus cedendi ictui}
Quia non est summa ullius corporis liquiditas,
sed cuidam consistentiae%
\indexS{liquiditas mixta consistentiae}%
\indexS{consistentia mixta liquiditati}
\edlabel{LH_37_04_033_033r_lgl-1}%
mixta.%
\edtext{%
}{%
\xxref{LH_37_04_033_033r_lgl-1}{LH_37_04_033_033r_lgl-2}%
\lemma{mixta.}%
\Bfootnote{%
\bracks{/} %
\textit{(1)}~Un\angles{a} superest %
\textit{(2)}~Ideo %
\textit{(a)}~ictus %
\textit{(b)}~luminis actio ita excipienda est,%
~\textit{L}%
}}
\pend%
\pstart%
Ideo luminis actio%
\indexS{actio luminis}%
\indexS{lumen agens}
ita excipienda est%
\edlabel{LH_37_04_033_033r_lgl-2}
ac si densissimum%
\bracks{,}
ac si vis aheneum%
\bracks{,}
murum%
\indexS{murus aheneus}
malleo percuterem%
\indexS{malleus}%
\bracks{:}
certum est
ictum recta continuatum iri,%
\indexS{ictus continuatus}
corpore quantaecunque sit duritiei,%
\indexS{durities corporis}
modo non infinitae,%
\indexS{durities infinita}
nonnihil cessuro.%
\indexS{corpus cedens ictui}
Si enim corpus sit \textit{efgh} planum,
ictusque incidens sit \textit{ab},%
\indexS{ictus incidens}
conabitur certe plicare corpus%
\indexS{corpus plicatum}
in statum \textit{ecfgihe},%
\indexS{status corporis}
quodque conabitur,
efficiet primo certe momento.%
\indexS{momentum primum}
Eodem momento autem corpus reaget,%
\indexS{corpus reagens}
seu se restituet,%
\indexS{corpus restituens se}
ac proinde omnia intra conatus transigentur.%
\indexS{conatus plicandi}%
\indexS{conatus restituendi se}
Manifestum est quoque
nihil ad rem pertinere corporis percussi densitatem,%
\indexS{densitas corporis percussi}%
\indexS{corpus percussum}
modo purum homogeneumque sit,%
\indexS{corpus purum}%
\indexS{corpus homogeneum}
ne ob varias superficies%
\indexS{superficies varia}
in %
alias lineas ictus%
\indexS{ictus discedens}
\edtext{%
discedat.
Semper enim is tantum ictus incidere%
\indexS{ictus incidens}
putandus est,
qui est ad superficiem obstaculi perpendicularis.%
\indexS{ictus perpendicularis ad obstaculum}%
\indexS{superficies obstaculi}%
\indexS{obstaculum}
Ideo in omni refractione ac ref\bracks{lex}ione%
\indexS{refractio}%
\indexS{reflexio}
sequeretur ictus partem perdi,%
\indexS{pars ictus perdita}
scilicet motum plano obstaculi parallelum,%
\indexS{motus parallelus}%
\indexS{obstaculum}
cum tamen constet
eum ab altero latere superesse,
et in citeriore non sentiri,
alioqui si in plano perpendiculari pergeret,
alio quodam plano opposito sentiri posset.
Necesse est ergo
ictum conatumve transire%
\indexS{ictus transiens in linea incidentiae}%
\indexS{conatus transiens in linea incidentiae}
in ea linea
qua incidit.%
}{%
\lemma{discedat.}%
\Bfootnote{%
\textit{(1)}~Potest %
\textit{(2)}~Debet tamen dici ictum illum minimum nunquam in alias lineas discedere, sed %
\textit{(3)}~Semper enim %
\bracks{...} obstaculi perpendicularis. %
\textit{(a)}~At resistentia %
\textit{(b)}~Ideo in omni refractione ac %
\bars{refractione \textit{ändert Hrsg.}} sequeretur ictus %
\bracks{...} qua incidit.%
~\textit{L}%
}}
Sed unde jam inclinatio ad
\edtext{%
perpendicularem,
vel a perpendiculari?%
}{%
\lemma{perpendicularem,}%
\Bfootnote{%
\textit{(1)}~restitutioque %
\textit{(2)}~vel a perpendiculari?%
~\textit{L}%
}}
\pend%
\begin{edfigure}%
\includegraphics%
{gesamttex/edit_VIII-4/images/LH_37_04_033_d_033r.pdf}%
\edcaption[0.250em]{%
\bracks{\textit{Fig.\:1}}%
}%
\end{edfigure}%
\pstart%
Si per Undas rem explicamus,%
\indexS{unda}
si materiam per poros diffusam credimus,%
\indexS{materia diffusa per poros}%
\indexS{pori}
non video
cur lumen eat in linea recta,%
\indexS{lumen iens in linea recta}
cur non in obliquis innumeris,
ut sonus.%
\indexS{sonus}
Nam si
fortissimum esse ictum%
\indexS{ictus fortissimus}
dicas,
deflecti impatientem,%
\indexS{ictus impatiens deflecti}
utique aut recta transit,%
\indexS{ictus transiens recta}
aut obstaculo reperto reflectitur.%
\indexS{ictus reflexus obstaculo}
Si simplex ille conatus propagandi motum%
\indexS{conatus simplex}%
\indexS{conatus propagandi motum}%
\indexS{motus propagatus}
causa esset luminis,%
\indexS{causa luminis}
omnia pertransiret,
sed iste simplex %
conatus%
\indexS{conatus simplex}
cum sit momentaneus%
\indexS{conatus momentaneus}
est insensibilis.%
\indexS{conatus insensibilis}
Opus ergo plurium concursu,%
\indexS{concursus plurium conatuum}%
\indexS{conatus plurimus}%
\indexS{conatus concurrentes}
quoniam tempore opus ad sentiendum agendumque.
Continuis ergo conatibus%
\indexS{conatus continuus}
et reactionibus%
\indexS{reactio continua}
seu renisibus%
\indexS{renisus continuus}
ipsius superficiei \textit{hg},
fit ut conatus agentis%
\indexS{conatus agentis}
continue crescat,%
\indexS{conatus crescens continue}
ut hominis portam effringentis.%
\indexS{homo effringens portam}%
\indexS{porta}
Et probe notandum
non sufficere
motum esse continuo supplemento auctum,%
\indexS{motus auctus supplemento continuo}
ad hoc
ut ictus ex resistentia%
\indexS{ictus fortior e resistentia}
fiat fortior,
nam unus quisque conatuum separatim destruitur,%
\indexS{conatus destructus}
et ut actiones ita et reactiones repetuntur.%
\indexS{actio et reactio}%
\indexS{actio repetita}%
\indexS{reactio repetita}
Id potius necesse est,
ut id quod agit,%
\indexS{agens quod sentit resistentiam}
resistentiam sentiens,%
\indexS{resistentia agenti}
velut iratam vim intendat.%
\indexS{vis irata}
Quod unde fiat,
explicandum est.
\pend%
\newpage%
\pstart%
Hoc autem nulla re rectius explicari potest,
quam vulgari experimento%
\indexS{experimentum vulgare}%
\indexS{experimentum notum infantibus}
infantibus quoque noto%
\indexS{infans}%
\bracks{:}
canalis angusti utrumque latus%
\indexS{canalis obstructus}
tenaci quadam%
\indexS{materia tenax}
\edtext{%
materia compacta,%
\indexS{materia compacta}%
}{%
\lemma{materia}%
\Bfootnote{%
\bars{, ut charta,\indexS{charta} \textit{gestr.}} compacta, %
~\textit{L}%
}}
mobili tamen,%
\indexS{materia mobilis}
si vim adhibeas%
\indexS{vis adhibita}
\edtext{%
obstru\bracks{itur}.%
\indexS{canalis obstructus}%
}{%
\lemma{obstruunt.}%
\Bfootnote{%
\textit{L\;ändert Hrsg.}%
}}
Inde embolo immisso,%
\indexS{embolus immissus}
alterum \textls{obstrumentum}%
\indexS{obstrumentum}
versus alterum
\edtext{%
ag\bracks{i}t,%
}{%
\lemma{agunt,}%
\Bfootnote{%
\textit{L\;ändert Hrsg.}%
}}
tamdiu
donec aer interceptus%
\indexS{aer interceptus}
auctae continue condensationis,%
\indexS{condensatio aucta continue}
obstrumenti adhaesionem vincat,%
\indexS{adhaesio obstrumenti}%
\indexS{obstrumentum}
majoreque id longe vi%
\indexS{vis aeris interceptis}
expellat,
quam si ab initio statim superavisset.
Quoniam hoc modo omnes nisus simul conservantur,%
\indexS{nisus conservatus}
quod repetitis simpliciter ictibus non fit.%
\indexS{ictus repetitus}
Eandem causam similemve esse
necesse est,
cur lumen ab obstaculo%
\indexS{obstaculum}
capiat vires,%
\indexS{lumen capiens vires ab obstaculo}
et quod resistit,
fortius plicet.
\pend%
\vspace{3.0em}%
\pstart%
\noindent%
\centering%
*     *     *
\pend%
\vspace{3.0em}%
\selectlanguage{german}%
\frenchspacing%
\footnotesize%
\pstart%
\noindent%
\textls{Ausführliche Datierungsgründe:}%
\pagelabel{Cc2-0481_dat}
Der \textit{Catalogue critique} (N.\;481) gibt als Datierung der vorliegenden Aufzeichnung N.\;\ref{Cc2-0481} \textit{vor August 1673} an.
Der Textträger weist ein Wasserzeichen auf, welches sich eindeutig in Leibnizens Pariser Zeit (1672--1676) einordnen lässt. %
Dieses Wasserzeichen kommt in verschiedenen Varianten vor, die insgesamt nur bis Anfang 1675 belegt sein dürften.
Auch inhaltlich weist die Aufzeichnung N.\;\ref{Cc2-0481} Ähnlichkeit mit anderen Texten aus der Pariser Zeit auf wie \textit{LSB} VIII,\;2 %
N.\;50 (letzte Monate 1674)\cite{01292} oder N.\;30 (Anfang 1675)\cite{01359}.
Die in N.\;\ref{Cc2-0481} aufgestellte Annahme, dass dichtere Körper, die einem Stoß langsamer nachgäben, einen stärkeren Impuls aufnehmen müssten (\refpassage{LH_37_04_033_033r_Zitate-2a}{LH_37_04_033_033r_Zitate-2b}) -- eine Anwendung der "allgemeinen Regel", dass die zusammengeballte Wirkung einer stetig zunehmenden Kraft (\textit{virtus}) desto stärker ausfalle, je länger ihr widerstanden werde (\refpassage{LH_37_04_033_033r_Zitate-1a}{LH_37_04_033_033r_Zitate-1b})~--, findet nämlich Anklang in VIII,\;2 N.\;50, S.\;441.22--23\cite{01292} (\textit{omne corpus pro ratione magnitudinis suae motui resistit, medii causa}) und wird in VIII,\;2 N.\;30, S.\;246--249\cite{01359} erörtert.
Die in der vorliegenden Aufzeichnung N.\;\ref{Cc2-0481} geäußerte Meinung, dass auch im luftleeren Raume eine bleierne Kugel weiter entfernt geschleudert werde als eine Feder  (\refpassage{LH_37_04_033_033r_Zitate-3a}{LH_37_04_033_033r_Zitate-3b}), lässt jedoch die dissipative Wirkung des widerstehenden Mediums außer Acht.
Dies dürfte wiederum als Zeichen gedeutet werden, dass die Aufzeichnung N.\;\ref{Cc2-0481} noch vor Leibnizens umfassender Auseinandersetzung mit diesem Phänomen im Jahre 1675 entstand (siehe \textit{LSB} VIII,\;2 N.\;30 bis N.\;38).\cite{01359}\cite{01360}\cite{01361}\cite{01362}\cite{01345}\cite{01364}\cite{01365}\cite{01366}\cite{01367}\cite{01368}\cite{01369}\cite{01370}\cite{01371}\cite{01372}\cite{01373}
Diese Betrachtungen lassen die Vermutung zu, dass die vorliegende Aufzeichnung N.\;\ref{Cc2-0481} frühestens noch im Jahr 1673 und spätestens im Jahr 1674 verfasst wurde.%
\pend%
\selectlanguage{latin}%
\frenchspacing%
\normalsize%
